\documentclass{article}

\usepackage[a4paper]{geometry}
\usepackage[english]{babel}
\usepackage{parskip}

\usepackage{amsmath, amssymb}
\usepackage{graphicx}
\usepackage{subcaption}
\usepackage{float}
\usepackage{hyperref}

\usepackage{listings}
\usepackage{xcolor}

\definecolor{codegreen}{rgb}{0,0.6,0}
\definecolor{codegray}{rgb}{0.5,0.5,0.5}
\definecolor{codepurple}{rgb}{0.58,0,0.82}
\definecolor{backcolour}{rgb}{0.95,0.95,0.92}

\lstdefinestyle{mystyle}{
  backgroundcolor=\color{backcolour},
  commentstyle=\color{codegreen},
  keywordstyle=\color{magenta},
  numberstyle=\tiny\color{codegray},
  stringstyle=\color{codepurple},
  basicstyle=\ttfamily\footnotesize,
  breaklines=true,
  captionpos=b,
  keepspaces=true,
  numbers=left,
  numbersep=5pt,
  showstringspaces=false,
  tabsize=2
}
\lstset{style=mystyle}

\newcommand{\PlotDir}{Plots}

\title{NURb Hand-in 1}
\author{Niels de Klerk (s3640477)}
\date{22-6-2026}

\begin{document}
\maketitle

\section{Simulating the solar system}

\subsection{(1a) }

Using \textit{astropy}, we find the current positions and velocities of the 8 planets and the Sun. We plot the positions of the planets and the Sun in both the x-y plane and the x-z plane below.

\begin{figure}[H]
    \centering
    \includegraphics[width=0.85\textwidth]{\PlotDir/initial_positions.png}
    \caption{Initial positions of the objects in the solar system in the x-y and x-z planes.}
\end{figure}

\subsection{(1b) }

We want to simulate the solar system 300 years in the future. For this, we will use the Leapfrog method as this is a symplectic integrator, which means that it approximately conserves the energy of a system, unlike methods such as Euler's method, RK4, etc. So, leapfrog is a good choice to simulate physical systems such as orbits. Leapfrog offsets the velocity and position in the steps in time. It starts by 'kicking' the velocity half a timestep forward. 
$$v_{i+1/2}=v_i+0.5dt*a_i$$
Leapfrog updates the position and velocity alternately using the previous value of the other.
$$v_{i+1/2}=v_{i-1/2}+dt*a_i$$
$$x_{i+1}=x_{i}+dt*v_{i+1/2}$$
This is repeated until the desired time/ number of steps is reached.
For the kick, we use RK4, which is explained further in 1c. For the calculation of the acceleration, the equation 
$$\textbf{f} = -G * \textbf{r} / (|r|^3)$$
is used, which gives the mass-normalized force. We loop over the bodies, and we calculate this force from this body to all the next ones. Then the acceleration due to all the next bodies is added to the total acceleration of this body by multiplying $\textbf{f}$ with the corresponding mass, and the acceleration due to this body is added to all the next bodies' acceleration. This is repeated for all bodies, only with respect to the next bodies in the list, so as not to double-count any force.
Both the orbits in the x-y plane, as well as the z position over time, are shown below. 

\begin{figure}[H]
    \centering
    \includegraphics[width=0.85\textwidth]{\PlotDir/orbits_xy_leapfrog.png}
    \caption{Orbits of the objects in the x-y plane using the leapfrog integrator.}
\end{figure}

The orbits from the simulations look to be correct and nicely closed. The only thing that stands out is that the orbit of Mercury seems to be less stable than the other orbits. This is probably due to Mercury being close to the Sun, so the forces acting on this planet are stronger. Furthermore, the orbit is more strongly curved than, for example, the orbit of Neptune. To resolve this problem, you would need to take smaller time steps to make the result more accurate.
 
\begin{figure}[H]
    \centering
    \includegraphics[width=0.85\textwidth]{\PlotDir/z_vs_time_leapfrog.png}
    \caption{Position in the z plane as a function of time for all objects using the leapfrog integrator.}
\end{figure}

From the simulations, we also find that the z-position of the planets and the Sun, i.e., the height above the midplane, oscillates. This is expected as the planets follow elliptic orbits at an angle with the midplane, which results in these oscillations.  

\subsection{(1c) }

Explain the integration algorithm you used. What are the differences compared to the leapfrog method? 

Make the same plots that you did in (b), as well as a plot of the difference in the bodies' x positions versus time between the leapfrog method and the other method. Discuss your results.

We want to compare the results we got from the leapfrog method to a different integrator to see the difference between the two. We will use RK4 as this is a commonly used integrator. In the case of the orbit, we use 
$$\textbf{y}=\begin{pmatrix}
\textbf{x} \\
\textbf{v}
\end{pmatrix}, f(\textbf{x}, \textbf{v})=\begin{pmatrix}
\textbf{v} \\
\textbf{a}(\textbf{x})
\end{pmatrix}$$
where $\textbf{a}(\textbf{x)}$ is the acceleration at $\textbf{x}$ due to all the other objects. 
so we can write the equations of motion as 
$$\frac{\mathrm{d}\textbf{y}}{\mathrm{dt}} = f(\textbf{x}, \textbf{v})$$
$\textbf{y}$ is then updated using RK4 be taking a weighted sum of the function at 4 points following 
$$k_1 = \mathrm{dt} * f(\textbf{y}_i)$$
$$k_2 = \mathrm{dt} * f(\textbf{y}_i+\frac{1}{2}k_1)$$
$$k_3 = \mathrm{dt} * f(\textbf{y}_i+\frac{1}{2}k_2)$$
$$k_4 = \mathrm{dt} * f(\textbf{y}_i+k_3)$$
The new position and velocity is then given by
$$\textbf{y}_{i+1} = \textbf{y}_i + (k_1 +2k_2 + 2k_3 + k_4)/6$$

The results of the RK4 simulation are shown below. 

\begin{figure}[H]
    \centering
    \includegraphics[width=0.8\textwidth]{\PlotDir/orbits_xy_another_method.png}
    \caption{Orbits of the objects in the x-y plane using another integration method.}
\end{figure}

The orbits from the simulations look to be correct and nicely closed. Contrary to the simulation from the leapfrog method, the orbit of Mercury seems to be a lot more stable in the RK4. This is probably due to the fact that RK4 is 4th-order accurate, whereas leapfrog is only 2nd-order. 

\begin{figure}[H]
    \centering
    \includegraphics[width=0.8\textwidth]{\PlotDir/z_vs_time_another_method.png}
    \caption{Position in the z plane as a function of time for all objects using another integration method.}
\end{figure}

Again, the z-position of the objects oscillates through the midplane of the solar system. However, there is no visual difference between these results and the results derived from the leapfrog method.

To directly compare the 2 methods, we calculate the difference in the x-position between the 2 methods. 

\begin{figure}[H]
    \centering
    \includegraphics[width=0.8\textwidth]{\PlotDir/x_difference_another_method_minus_leapfrog.png}
    \caption{Difference in the x positions of the bodies as a function of time between the leapfrog method and another integration method.}
\end{figure}

Here we see that the difference between the methods slowly increases over time. Notably, the differences seem to oscillate. 

\subsection{(1d) }
We have also made a video of the simulation. This can be found in the folder with all the plots. 

\section{Calculating potentials in different ways}

\subsection{(2a) }

We have a box of size 250 Mpc with 256$^3$ particles. We want to construct an octree to partition these particles into nodes. The construction of the tree starts with the root node that contains all of the particles with index (0,0,0). Next, we split each direction into 2, resulting in 8 octants. For each octant, particles are assigned to it, and the total mass and center of mass are calculated. To keep track of the nodes, they are indexed by using the index of the parent node (x$_i$, y$_i$, z$_i$) and give this node the index (2x$_i$+x$_\mathrm{offset}$, 2y$_i$+y$_\mathrm{offset}$, 2z$_i$+z$_\mathrm{offset}$), where x$_\mathrm{offset}$ is 0 if the node is left of the center of the parent node and 1 otherwise, similarly for the other 2 axes. When looking at the binary representation of the index, for example, in the x-axis, this gives a series of left/right steps in the x-axis that correspond to the node. The splitting of the nodes is repeated until a node is empty or a maximum depth is reached. 

Using the indices of the nodes, we can directly find the mass in certain regions in the box. Using this, we plot different slices at levels 2, 5, and 7.  

\begin{figure}[H]
    \centering
    \includegraphics[width=0.8\textwidth]{\PlotDir/fig2a_level3.png}
    \caption{Mass distribution for the first four slices in xi at level 3 of the octree.}
\end{figure}

In the slices at level 3, there is little to no structure visible, as this resolution is too low to resolve any. 

\begin{figure}[H]
    \centering
    \includegraphics[width=0.8\textwidth]{\PlotDir/fig2a_level5.png}
    \caption{Mass distribution for the first four slices in xi at level 5 of the octree.}
\end{figure}

In the slices at level 5, there is a little more structure than for the previous slices, but the structure is still not very clear.
\begin{figure}[H]
    \centering
    \includegraphics[width=0.8\textwidth]{\PlotDir/fig2a_level7.png}
    \caption{Mass distribution for the first four slices in xi at level 7 of the octree.}
\end{figure}

In the slices at level 7, the filamentary structure of the dark matter is now visible.

\subsection{(2b) }
We can calculate the potential that the particles create by using
$$\Phi=-\frac{G}{\pi}\mathrm{iFT}[\mathrm{FT}[\rho]/k^2]$$
So we need the Fourier Transform (FT) to find the potential. 
For the FT, we use the Cooley-Tukey algorithm. 
First, we use bit reversal to swap the elements in the array. Next, we combine the elements using different complex weights until we recover the FT of the array.
The inverse FT can be found using the FT
$$\mathrm{iFT}[f]=\frac{\mathrm{FT}(f^*)^*}{m},$$
where $^*$ denotes conjugation and $m$ is the number of points in $f$.
To apply both the FT and iFT in 3 dimensions, the FT/iFT is applied along each dimension in sequence.
To find $\rho$, we use the bottom level of the octree, which subdivides the space into a 128$^3$ grid of points. So the mass of each node is assigned to the grid points, which gives a density grid. Next, we find the wavenumber of the different modes in Fourier space. and divide the FT of the density grid by 
$$k^2=k_x^2+k_y^2+k_z^2.$$
This result holds for all grid points except at the point where the wavenumber in each direction is 0, which would divide the FT of the density field by 0. However, we can just set this mode to 0, as the potential is defined up to a scalar, which gives the freedom to choose this one mode. Next, we apply the iFT to the rescaled modes and multiply the result by $-\frac{G}{\pi}$, which gives the potential. 4 slices of the potential are shown below.
  
\begin{figure}[H]
    \centering
    \includegraphics[width=0.8\textwidth]{\PlotDir/fig2b.png}
    \caption{Potential computed using the FFT method for different slices.}
\end{figure}

The potential is much smoother than the filamentary structure of the dark matter.

\section{Spiral and elliptical galaxies}

\subsection{(3a) }

We want to classify galaxies into 2 different groups: spiral and elliptical galaxies. For this, we use logistic regression for binary classification. Before we train the model, we feature-scale the data using standardization. 
$$x\rightarrow \frac{x-\mu}{\sigma},$$
where $\mu$ is the mean of the values and $\sigma$ the standard deviation of the values.
The results of the standardization are shown in the figure below. 

\begin{figure}[H]
    \centering
    \includegraphics[width=0.8\textwidth]{\PlotDir/fig3a.png}
    \caption{Distributions of the rescaled features for the spiral and elliptical galaxies.}
\end{figure}

In the distributions, we see that the extended feature and the emission line flux are sharply peaked. This is because there are a small number of points that are much further away from the mean, which increases the variance and thus makes the rescaled feature more sharply peaked.
\lstinputlisting[
        firstline=1,
        lastline=11,
        basicstyle=\ttfamily\small,
        frame=single
    ]{\PlotDir/galaxy_data_scaled.txt}

\subsection{(3b) }

Explain how you implemented logistic regression. 
Plot the evolution of your cost function with the number of iterations for different pairs of features.
Explain how does it change depending on the feature pair.

To apply the logistic regression we use an linear model that performs a weight sum of the features and adds a bias to the result.
$$z= \textbf{$\theta$}^T\textbf{x}+\beta,$$
where $\textbf{$\theta$}$ are the weight, $\textbf{x}$ the features and $\beta$ the bias.
The result is then put through the sigmoid activation function 
$$\hat{y}=\frac{1}{1+\exp z},$$
where $\hat{y}$ is used as the prediction. To improve the model, we want to minimize the logistic loss function 
$$J = -\frac{1}{m}\sum_{i=0}^{m-1}y_i\ln{\hat{y}_i}+(1-y_i)\ln{(1-\hat{y}_i)}$$
Downhill simplex is used to minimize the loss function, and the resulting evolution of the loss for each iteration is plotted below.

\begin{figure}[H]
    \centering
    \includegraphics[width=0.8\textwidth]{\PlotDir/fig3b.png}
    \caption{Evolution of the cost function with the number of iterations for different pairs of features in the logistic regression model.}
\end{figure}

Here we see that, depending on the combination of features used, the loss levels off at different values. So from this we see that the feature combination 1+2 ($\kappa_\mathrm{CO}$ + Color) gives a much better split than, for example, feature combination 3+4 (Extended+Emission line flux). 

\subsection{(3c) }

Building the confusion matrix for the different features combinations gives
\begin{table}[H]
\centering
\caption{The true/false positive/negative results and the F1-score for the different feature combinations}
\begin{tabular}{|c|c|c|c|c|c|}
\hline
Feature Combination & TP & FN & TN & FN & F1-score \\
\hline
\input{Plots/logistic_regression_metrics0.txt} \\
\input{Plots/logistic_regression_metrics1.txt} \\
\input{Plots/logistic_regression_metrics2.txt} \\
\input{Plots/logistic_regression_metrics3.txt} \\
\input{Plots/logistic_regression_metrics4.txt} \\
\input{Plots/logistic_regression_metrics5.txt} \\
\hline
\end{tabular}   
\end{table}
And plotting the different feature  combinations with the decision boundary gives
\begin{figure}[H]
    \centering
    \includegraphics[width=0.7\textwidth]{\PlotDir/fig3c.png}
    \caption{Plot of the pairs of two features against each other with the decision boundary.}
\end{figure}
From both the table and the figure, $\kappa_\mathrm{CO}$ + Color gives the best split, because $\kappa_\mathrm{CO}$ itself splits the data good into spiral and galaxy type. Again, the emission line flux and extended combination give the worst results. Visually, this can be explained by the fact that the elliptical galaxies span across this entire range, while the spiral galaxies are more centered. So a linear model would not give the best results.


\end{document}